\section{Operator interpretation}
\label{app:operator}

Define $\mathcal A:\V\to\V'$ by
\begin{equation}
  \dual{\mathcal Aw}{v}_{\V',\V}=a(w,v).
  \label{eq:operator-definition}
\end{equation}
Continuity implies $\norm{\mathcal A}_{\mathcal L(\V,\V')}\leq M$, while
coercivity and the Lax--Milgram theorem imply that $\mathcal A$ is invertible
with $\norm{\mathcal A^{-1}}\leq m^{-1}$.

Let $P_h:\V\to\Vh$ be the Ritz projection defined by
\begin{equation}
  a(P_hw,v_h)=a(w,v_h)
  \qquad\forall v_h\in\Vh.
  \label{eq:ritz-projection}
\end{equation}
Then $u_h=P_hu$, and C\'ea's lemma gives
\begin{equation}
  \norm{I-P_h}_{\mathcal L(\V,\V)}
  \leq \frac{M}{m}.
  \label{eq:ritz-bound}
\end{equation}
For a symmetric bilinear form, $P_h$ is orthogonal in the energy inner
product, so the sharper best-approximation identity
\begin{equation}
  \energy{u-P_hu}
  =\inf_{v_h\in\Vh}\energy{u-v_h}
  \label{eq:energy-best-approximation}
\end{equation}
holds.
